\documentclass[11pt]{article}
\usepackage[T1]{fontenc}
\usepackage{lmodern,amsmath,amssymb,amsthm,booktabs}
\usepackage[margin=1in]{geometry}
\usepackage[colorlinks=true,urlcolor=blue,citecolor=blue,linkcolor=blue]{hyperref}
\newtheorem{theorem}{Theorem}
\newtheorem{corollary}[theorem]{Corollary}
\newcommand{\Q}{\mathbb Q}
\newcommand{\N}{\mathbb N}
\title{Barry's shifted Hankel conjecture and an extension of Cigler's formula}
\author{Haobo Ma\\\small ChronoAI Pte Ltd, Singapore
\and Wenlin Zhang\\\small National University of Singapore\\\small ChronoAI Pte Ltd, Singapore}
\date{November 16, 2026}
\begin{document}
\maketitle
\begin{abstract}
We prove Paul Barry's shifted Hankel conjecture in OEIS A188475,
revision 2 (April 1, 2011): for every $n\geq0$, the determinant of
the recurrence-defined matrix $[q(i+j+2)]_{0\leq i,j\leq n}$ equals
$(2n^3+3n^2+n+3)/3$, by a moment-matrix cofactor argument extending
Cigler's unit-descent formula (2011, Theorem~A*).
The recurrence gives a tridiagonal moment model with squared norms
$1,2,2,\ldots$; its monic polynomials have explicitly computable constant
terms. The same calculation gives both adjacent Hankel shifts and a
one-parameter determinant formula. A direct bijection connects the
recurrence with the source's nonexchangeable subsequences and Callan's
restricted Grand Dyck paths.
\end{abstract}

\section{The source and the determinant}
Barry's comment in A188475 states ``Hankel transform of A137398(n+1)
(conjecture)'' \cite{barry}. The latter entry changed its offset
from zero to one on June 7, 2021 \cite{source}. Consequently the
original shifted assertion is about $q(i+j+2)$ in the current numbering,
not $q(i+j+1)$. Define $C_k=\frac1{k+1}\binom{2k}{k}$ and set
\begin{equation}\label{rec}
\begin{split}
q(0)&=q(1)=0,\qquad q(2)=1,\qquad q(3)=2,\\
q(r)&=2q(r-1)+2q(r-2)+\sum_{k=1}^{r-3}C_kq(r-k-1)
\quad(r\geq4).
\end{split}
\end{equation}
The auxiliary $q(0)$ will be used only for an adjacent-shift identity.
The first positive-index terms are $0,1,2,7,22,74,252$.

\begin{theorem}\label{main}
For every integer $N\geq0$, with the determinant of an empty matrix
equal to one,
\begin{equation}\label{sizeN}
\det[q(i+j+2)]_{0\leq i,j<N}
 =1+2\sum_{k=0}^{N-1}k^2.
\end{equation}
In particular, for every $n\in\N$,
\begin{equation}\label{cubic}
\det[q(i+j+2)]_{0\leq i,j\leq n}
 =1+2\sum_{k=0}^{n}k^2
 =\frac{2n^3+3n^2+n+3}{3}.
\end{equation}
\end{theorem}

The key is to adjoin a constant moment to twice the source generating
function. This gives a particularly simple polynomial recurrence.
Goy and Shattuck \cite[Theorem 3, p.~425]{goy} express A137398
as Hessenberg--Toeplitz determinants of Fine numbers; those matrices
are different from the Hankel matrix in \eqref{sizeN}. Here a
moment-matrix cofactor evaluates the latter directly. The proof below
derives the required orthogonality algebraically, without a measure
representation or a nonsingularity assumption on the source matrices.

The unit-descent case is already contained in Cigler's
Theorem~A*, equations (1.15)--(1.18) \cite{cigler}, and in
Chang--Hu--Zhang's Theorem~3.5 \cite{chang}. In their notation,
take $(a,b,t)=(2,1,-2)$, so their tail is $T=C^2$ and their
full moment series is $(1-z^2T)^{-1}$. Their shift-two formula
then gives $1+\sum_{k=0}^{N-1}k^2$.
Barry's recurrence instead requires the initial descent weight two,
giving $1+2\sum_{k=0}^{N-1}k^2$. Corollary~\ref{parameter}
includes both cases and proves the formula for every initial descent
weight $\lambda$.

Closely related work by Mu and Wang \cite{mu-wang} concerns Hankel
determinants of shifted Catalan-like numbers. Its relationship to the cubic
identity \eqref{cubic} remains unresolved; overlap is not excluded.
No priority claim is made.

\section{A moment model from the literal recurrence}
All power series in this section lie in $\Q[[z]]$. Write
\[
C(z)=\sum_{k\geq0}C_kz^k,\quad T(z)=C(z)^2,\quad
F(z)=\sum_{r\geq0}q(r)z^r.
\]
The Catalan identity $C=1+zC^2$ implies
$T=(1+zT)^2$, and hence
\begin{equation}\label{tail}
T(1-2z-z^2T)=1.
\end{equation}
Summing \eqref{rec}, including its initial values, gives
\begin{equation}\label{Feq}
F=z^2+2zF+2z^2F+z(C-1)F
 =z^2+2zF+2z^2F+z^2TF.
\end{equation}
All denominators below have constant term one. Using \eqref{tail}
in \eqref{Feq} yields
\begin{equation}\label{moment}
F=\frac{z^2T}{1-2z^2T},\qquad
M:=1+2F=\frac1{1-2z^2T}.
\end{equation}
Thus $m_r:=[z^r]M$ satisfies $m_0=1$ and $m_r=2q(r)$ for $r\geq1$.

For completeness, construct its polynomial model on the vector space
$V=\bigoplus_{j\geq0}\Q e_j$. Define
\[
Je_0=e_1,\qquad
Je_j=e_{j+1}+2e_j+b_je_{j-1}\quad(j\geq1),\qquad
b_1=2,\quad b_j=1\ (j\geq2).
\]
The coefficient of $e_0$ in $J^re_0$ is the sum of the weights of
length-$r$ paths on the nonnegative integers. An upward step has
weight one, a horizontal step at positive height has weight two,
and a downward step from height $j$ has weight $b_j$; there is no
horizontal step at zero. Excursions above height zero, shifted down
by one, have generating function $S$ satisfying
$S=1+2zS+z^2S^2$: the nonempty path begins either with a horizontal
step or with an upstep and its matching downstep. This equation
has a unique constant-term-one solution, determined coefficient by
coefficient, so $S=T$ by \eqref{tail}. Every path from zero to zero
is a concatenation of excursions of weight $2z^2T$. Its generating
function is therefore $M$ in \eqref{moment}. In particular,
\begin{equation}\label{L}
L(x^r):=m_r=[e_0]J^re_0
\end{equation}
defines a linear functional $L:\Q[x]\to\Q$.

Give $V$ the symmetric bilinear form
$\langle e_i,e_j\rangle=h_i\delta_{ij}$ with $h_0=1$ and
$h_j=2$ for $j\geq1$. The identity $h_{j+1}=b_{j+1}h_j$
shows on each pair of basis vectors that $J$ is self-adjoint for
this form. Define monic polynomials by
\begin{equation}\label{P}
P_0=1,\quad P_1=x,\quad
P_{j+1}=(x-2)P_j-b_jP_{j-1}\quad(j\geq1).
\end{equation}
Induction gives $P_j(J)e_0=e_j$: the cases $j=0,1$ follow from
the definitions, and the three terms in \eqref{P} cancel the
$e_j$ and $e_{j-1}$ components at the induction step. Consequently
\begin{equation}\label{orth}
L(P_iP_j)=\langle e_0,P_i(J)P_j(J)e_0\rangle
 =\langle e_i,e_j\rangle=h_i\delta_{ij}.
\end{equation}
Every vector and operator expression used here has finite support.

\section{The cofactor evaluation}
Put $G_{N+1}=[m_{i+j}]_{0\leq i,j\leq N}$, and let
$A_{ij}=[x^i]P_j$ on the same indices. Monicity makes $A$ upper
triangular with diagonal one. Equation \eqref{orth} gives
\[
A^{\mathsf T}G_{N+1}A=D,\qquad
D=\operatorname{diag}(1,2,\ldots,2).
\]
It follows that $\det G_{N+1}=2^N$ and
$G_{N+1}^{-1}=AD^{-1}A^{\mathsf T}$. In particular the
$(0,0)$ cofactor is
\begin{equation}\label{cofactor}
\operatorname{cof}_{00}(G_{N+1})
 =2^N\sum_{j=0}^N\frac{P_j(0)^2}{h_j}.
\end{equation}
The polynomial recurrence gives $P_1(0)=0$ and $P_2(0)=-2$.
For $j\geq2$ it gives $P_{j+1}(0)=-2P_j(0)-P_{j-1}(0)$.
These initial values and induction imply
\begin{equation}\label{constant}
P_{k+1}(0)=2(-1)^k k\quad(k\geq0).
\end{equation}
Indeed substituting the formulas for two consecutive terms in
the last recurrence gives the formula for the next term.
Thus the right side of \eqref{cofactor} is
\[
2^N\left(1+2\sum_{k=0}^{N-1}k^2\right).
\]
Deleting row and column zero of $G_{N+1}$ leaves
$[m_{i+j+2}]_{0\leq i,j<N}=2[q(i+j+2)]_{0\leq i,j<N}$.
Its determinant is $2^N$ times the determinant in \eqref{sizeN}.
Canceling $2^N$ proves that identity, including $N=0$.
Setting $N=n+1$ and using
$\sum_{k=0}^n k^2=n(n+1)(2n+1)/6$ proves \eqref{cubic}.
The sum-of-squares formula follows by induction, since its value
increases by $(n+1)^2$ when $n$ is replaced by $n+1$.

\section{Adjacent shifts and a parameter extension}
\begin{corollary}\label{adjacent}
For every $N\geq1$,
\begin{align}
\det[q(i+j+1)]_{0\leq i,j<N}&=1-N,\label{shift1}\\
\det[q(i+j)]_{0\leq i,j<N}&=-\sum_{k=0}^{N-2}k^2,\label{shift0}
\end{align}
where the last sum is empty for $N=1$.
\end{corollary}
\begin{proof}
Let $A_N$ be the polynomial coefficient matrix on indices $0,\ldots,N-1$
and let $J_N$ be the principal $N$ by $N$ block of $J$.
Multiplication by $x$ in the $P_j$ basis has these tridiagonal
coefficients. Thus \eqref{orth} implies
\[
A_N^{\mathsf T}[m_{i+j+1}]_{i,j<N}A_N
 =\operatorname{diag}(h_0,\ldots,h_{N-1})J_N.
\]
The term $P_N$ arising in $xP_{N-1}$ pairs to zero with every
$P_i$, $i<N$, so truncation introduces no extra term.
Expansion of $\det(xI-J_N)$ along the last row gives exactly
\eqref{P}, with initial determinants $1,x$; hence this determinant
is $P_N(x)$. It follows from \eqref{constant} that
\[
\det[m_{i+j+1}]_{i,j<N}
 =2^{N-1}(-1)^NP_N(0)=2^N(1-N).
\]
Dividing by $2^N$ proves \eqref{shift1}.
For \eqref{shift0}, the defining relation $M=1+2F$ gives
$2[q(i+j)]_{i,j<N}=G_N-E_{00}$, where $E_{00}$ has a single
one at $(0,0)$. Linearity in that entry yields
\[
2^N\det[q(i+j)]_{i,j<N}
 =\det G_N-\operatorname{cof}_{00}(G_N)
 =-2^N\sum_{k=0}^{N-2}k^2
\]
by \eqref{cofactor} with $N-1$ in place of $N$.
\end{proof}

\noindent In the original zero-based A137398 convention, its unshifted
Hankel transform is $-n$ for matrices of order $n+1$, while its
once-shifted transform is the cubic in \eqref{cubic}. This also
quantifies why applying the conjectural shift to the current offset
would give a different answer.

\begin{corollary}\label{parameter}
For $\lambda\in\mathbb C$, define $q_\lambda(r)$ by
\[
\sum_{r\geq0}q_\lambda(r)z^r
 =\frac{z^2C(z)^2}{1-\lambda z^2C(z)^2}.
\]
For every $N\geq1$,
\begin{align*}
\det[q_\lambda(i+j+2)]_{i,j<N}
 &=1+\lambda\sum_{k=0}^{N-1}k^2,\\
\det[q_\lambda(i+j+1)]_{i,j<N}&=1-N,\\
\det[q_\lambda(i+j)]_{i,j<N}&=-\sum_{k=0}^{N-2}k^2.
\end{align*}
In particular, $q_2=q$.
\end{corollary}
\begin{proof}
For $\lambda\ne0$, replace $b_1=2$ by $b_1=\lambda$ in the
polynomial model, leaving all other coefficients unchanged. The
moments become $M_\lambda=1+\lambda\sum q_\lambda(r)z^r$,
the squared norms become $1,\lambda,\lambda,\ldots$, and
$P_{k+1}(0)=\lambda(-1)^k k$. The same cofactor evaluation gives
$\lambda^N(1+\lambda\sum_{k<N}k^2)$, which is
$\lambda^N\det[q_\lambda(i+j+2)]_{i,j<N}$.
The two adjacent-shift calculations give respectively $1-N$ and
$-\sum_{k=0}^{N-2}k^2$ after cancellation. For fixed $N$, every
matrix entry and determinant is a polynomial in $\lambda$.
The identities for all nonzero $\lambda$ therefore also hold at zero.
\end{proof}

\section{The source's combinatorial interpretations}
\begin{corollary}\label{paths}
For $n\geq1$, $q(n)$ counts both the nonexchangeable pairs of
subsequences in the name of A137398 and the $2n$-step Grand Dyck paths
in Callan's November 25, 2021 comment: balanced up/down paths starting
with an upstep, with no peak at height $1$ and no valley at height $-1$.
Here, as in the source's examples, an admissible exchange must preserve
the order inherited from the ground sequence.
\end{corollary}
\begin{proof}
Relabel the ground set increasingly as $1,\ldots,2n$ and list each
subsequence increasingly; reversing both subsequences, if necessary,
preserves the exchange condition. Encode membership in the subsequence
containing $1$ by an upstep and membership in the other by a downstep.
This is a bijection with balanced paths starting up.

Write the two increasing subsequences as $p_1<\cdots<p_n$ and
$s_1<\cdots<s_n$. Swapping $p_i$ and $s_i$ preserves both orders
exactly when they are adjacent in the merged order: if an element lies
strictly between them, it belongs to one subsequence and obstructs its
order after the swap; adjacency conversely preserves every neighboring
inequality. Since both elements have rank $i$ in their own subsequence,
such an adjacent pair occupies positions $2i-1,2i$. It is therefore
precisely an up/down or down/up pair starting at height zero. Thus
nonexchangeability is equivalent to excluding peaks at $1$ and valleys
at $-1$.

For the enumeration, follow Callan's first-return decomposition
\cite{source}. A positive primitive excursion has the form $UAD$,
where $A$ is a nonempty Dyck path; the nonempty condition excludes its
peak at height $1$. Its generating function, with $z$ marking a pair
of steps, is $B=z(C-1)$. Negative primitive excursions are its
reflections and have the same series. The first excursion is positive,
and every later one has either sign. Hence the series of nonempty
allowed paths is $B/(1-2B)=z^2T/(1-2z^2T)=F$, proving the count.
\end{proof}

\section{Source coverage and generating functions}
\begin{center}\small
\begin{tabular}{@{}p{.43\linewidth}p{.47\linewidth}@{}}
\toprule
Source item & Coverage\\\midrule
A188475, Barry's shifted Hankel comment & Proved by Theorem~\ref{main}.\\
A137398, Catalan recurrence & Defines $q$; its moment identity is derived in
\eqref{Feq}--\eqref{moment}.\\
A137398, Barry's January 22, 2009 comment, ``g.f. appears to be g(x)-1''
& Refuted-and-corrected: in current indexing $F=(g-1)/2$, not $g-1$.\\
A188475, ordinary generating function & The cubic gives
$(1-z+5z^2-z^3)/(1-z)^4$.\\
A188475, Greubel's exponential generating function & Proved below from
the cubic.\\
A188475, order-four recurrence & Proved below by fourth differences.\\
A188475, relation to A006331 & The equality with one plus twice the
sum of squares follows from \eqref{cubic}.\\
A137398, radical generating function & Proved from
\eqref{Feq}--\eqref{moment}.\\
A137398, Callan's path count and equation & Reproved; prior proof
Callan (2021), reproduced in Corollary~\ref{paths}.\\
Adjacent determinants \eqref{shift1}--\eqref{shift0} & Not posed in these
entries; proved by Corollary~\ref{adjacent}.\\
A137398, nonexchangeable subsequences in the sequence name & Proved by
Corollary~\ref{paths}.\\
A137398, displayed triangle and its recurrences & Not covered.\\
Parameter $\lambda$ in Corollary~\ref{parameter} & Not posed in these
entries; proved for every complex $\lambda$.\\
Cigler's unit-descent specialization $\lambda=1$ & Reproved; prior
proof Cigler (2011), Theorem~A*, and Chang--Hu--Zhang (2013), Theorem~3.5.\\
Mu--Wang (2017), shifted Catalan-like determinants \cite{mu-wang}
& Closely related work; its relationship to the cubic identity
\eqref{cubic} remains unresolved. Overlap is not excluded.\\
\bottomrule
\end{tabular}
\end{center}
\noindent The ordinary generating-function row follows by applying
$z\frac{d}{dz}$ twice to $1/(1-z)$ and summing partial sums:
$\sum_{n\geq0}(1+2\sum_{k=0}^nk^2)z^n
=1/(1-z)+2z(1+z)/(1-z)^4$.
Writing $a(n)$ for the cubic, the exponential generating function is
\[
\sum_{n\geq0}a(n)\frac{z^n}{n!}
 =\frac{3+6z+9z^2+2z^3}{3}e^z.
\]
This follows by applying $(2\theta^3+3\theta^2+\theta+3)/3$
to $e^z$, where $\theta=z\frac{d}{dz}$. The fourth difference
of a cubic vanishes, so
$a(n)=4a(n-1)-6a(n-2)+4a(n-3)-a(n-4)$ for $n\geq4$.
Callan's equation is $F=z(C-1)+2z(C-1)F$, which is
\eqref{moment} because $C-1=zT$. Substituting
$z^2T=(1-2z-\sqrt{1-4z})/2$ in \eqref{Feq} gives the
source's radical expression
$F=2z^2/(1-2z-4z^2+\sqrt{1-4z})$.
For the continued-fraction row, Barry's $g$ has tail
$T=1/(1-2z-z^2T)$ and the full fraction is
$g=M=1/(1-2z^2T)$. Since $q(2)=1$ while $[z^2](g-1)=2$,
the factor two cannot be dropped.
In the original zero-based convention the series is $F/z=(g-1)/(2z)$.

\section*{Verification and availability}
The all-index statements and combinatorial interpretations have the
written proofs above. The accompanying \texttt{verify.py} checks exact
determinants, moment coefficients, polynomial orthogonality, parameter
specializations, and both combinatorial counts; \texttt{checks.json}
specifies the finite ranges.
The pinned Lean declaration
\href{https://github.com/the-omega-institute/trureturing/blob/c6d80ec1ad5920a0470d7ff831c17f6303e6e578/D5/S3/Combinatorics/CatalanPowerHankel/BarryA188475.lean}{\texttt{BarryA188475.result}}
states both equalities in \eqref{cubic} for the rational recurrence
\texttt{BarryA188475.q} and the matrix \texttt{BarryA188475.hankel}.
Its scope is Theorem~\ref{main} with $N=n+1$. The definitions
\texttt{q} and \texttt{hankel} specify the recurrence and shift;
\texttt{result} asserts their determinant identity for all $n$.
The adjacent shifts, parameter family and combinatorial corollary
are supplied by the written proofs, rather than by that declaration.

\begin{thebibliography}{9}
\small
\bibitem{barry} P. Barry, \emph{OEIS A188475}, April 1, 2011,
revision 2, comment and sequence name.
\url{https://oeis.org/A188475}.
Original comment: \url{https://oeis.org/history?seq=A188475&start=20}.
\bibitem{source} G. Beavers, G. Starling and W. Li,
\emph{OEIS A137398}, 2008, formulas and comments by P. Barry (2009)
and D. Callan (2021). \url{https://oeis.org/A137398}.
Offset correction, revision 11 (June 7, 2021):
\url{https://oeis.org/history?seq=A137398&start=20}.
\bibitem{goy} T. Goy and M. Shattuck,
\emph{Hessenberg--Toeplitz Matrix Determinants with Schr\"oder and Fine
Number Entries}, Carpathian Mathematical Publications
\textbf{15} (2023), no.~2, 420--436.
\url{https://doi.org/10.15330/cmp.15.2.420-436}.
\bibitem{squares} The OEIS Foundation,
\emph{OEIS A006331}, twice the sum of the first $n$ squares.
\url{https://oeis.org/A006331}.
\bibitem{cigler} J. Cigler, \emph{Some nice Hankel determinants},
2011, Theorem~A*, equations (1.15)--(1.18).
\url{https://arxiv.org/abs/1109.1449}.
\bibitem{chang} X.-K. Chang, X.-B. Hu and Y.-N. Zhang,
\emph{A direct method for evaluating some nice Hankel determinants and
proofs of several conjectures}, Linear Algebra and its Applications
\textbf{438} (2013), 2523--2541, Theorem~3.5.
\url{https://doi.org/10.1016/j.laa.2012.10.050}.
\bibitem{mu-wang} L. Mu and Y. Wang,
\emph{Hankel determinants of shifted Catalan-like numbers},
Discrete Mathematics \textbf{340} (2017), no.~6, 1389--1396.
\url{https://doi.org/10.1016/j.disc.2016.09.035}.
\end{thebibliography}
\medskip
\noindent\small\copyright\ 2026 Haobo Ma and Wenlin Zhang.
Licensed under Creative Commons Attribution 4.0 International (CC BY 4.0),
\url{https://creativecommons.org/licenses/by/4.0/}.
\end{document}
